\section{Fork/Lock Token Measurement}
\label{app:fork-lock-measurement}

For each model family, we measured fork- and lock-like token positions by
evaluating teacher next-token distributions on fixed student traces. We used the
same 60 OpenMathReasoning prompts for every model and generated one student
reasoning trace per prompt using the base student prompt. For each trace, we
evaluated the teacher distribution at every generated token position under three
conditioning settings: \emph{base}, \emph{sparse}, and \emph{dense}. The base
condition included only the problem statement; the sparse condition additionally
provided the correct final answer; and the dense condition provided privileged
context in the form of a truncated reference solution from a stronger model.

For each token position $t$, we formed the teacher context $(x, y_{<t})$ and
stored the teacher's retained top-$K$ next-token log-probabilities. In the final
SRT runs, we used $K=3$ and set \texttt{max\_student\_tokens} to 3072 for the
base and sparse conditions. For dense runs, we capped the reference trace at
2048 tokens and the student trace at 1536 tokens to avoid prompt-logprob memory
failures. All six cells per model were completed: OPSD and OPD crossed with
base, sparse, and dense, with 60 traces per cell.

\paragraph{Entropy-threshold analysis.}
We first normalized the entropy of the retained top-$K$ distribution:
\[
H_K^{\mathrm{norm}}
=
\frac{-\sum_{i=1}^{K} q_i \log q_i}{\log K},
\]
where $q_i$ denotes the top-$K$ probabilities renormalized over the retained
support. Positions with $H_K^{\mathrm{norm}} \le 0.20$ were labeled \emph{lock}
tokens, positions with $H_K^{\mathrm{norm}} \ge 0.60$ were labeled \emph{fork}
tokens, and all remaining positions were labeled neutral.

\paragraph{Support-aware SSD approximation.}
We also classified positions using the geometry of the retained support.
Starting from the saved top-$K$ distribution, we applied top-$p$ truncation with
$p=0.8$ and computed the retained support size, top-token probability,
top-1/top-2 log-probability gap, entropy-derived effective support size
\[
N_{\mathrm{eff}} = \exp(H_S),
\]
and the number of competitive tokens within a factor of 3 of the top token. A
position was labeled lock-like when the retained support was sharply
concentrated, and fork-like when multiple retained tokens remained competitive.
Tokens outside the retained support were treated as tail mass rather than forks.
Positions satisfying neither criterion were labeled neutral.

\paragraph{Aggregation.}
For each trace and conditioning setting, we computed fork, lock, and neutral
rates as the fraction of classified token positions in the trace. We visualize
per-trace rates using boxplots, separately for OPSD and OPD, with the base,
sparse, and dense conditions shown in each panel. Boxes summarize the
distribution across 60 traces; jittered points show individual traces; and
diamond markers indicate means.
